\chapter{Pengenalan Arsitektur Perangkat Lunak dan Client-Server}

\section*{Tujuan Pembelajaran}

\begin{enumerate}
  \item Menjelaskan arti arsitektur perangkat lunak beserta aspek yang
    membedakannya dari rancangan rinci.
  \item Membedakan tanggung jawab klien dari tanggung jawab server pada
    arsitektur client-server.
  \item Mengukur waktu bolak-balik satu permintaan dan perubahannya ketika
    jumlah klien bertambah.
\end{enumerate}

\section{Pendahuluan}

Setiap sistem perangkat lunak punya arsitektur, baik dirancang maupun tidak.
Arsitektur adalah susunan bagian-bagian besar sistem beserta aturan hubungan di
antaranya. Susunan itu menentukan hal-hal yang mahal diubah di kemudian hari,
misalnya batas antar modul dan cara modul berkomunikasi.

Yang membedakan arsitektur dari rancangan rinci adalah biaya perubahannya.
Mengganti nama sebuah fungsi murah. Memindahkan aturan bisnis dari satu bagian
ke bagian lain juga masih murah. Sebaliknya, memecah satu sistem menjadi
beberapa layanan terpisah menyentuh hampir seluruh berkas. Keputusan yang mahal
diubah itulah yang disebut keputusan arsitektural.

Bab ini membuka dengan pengertian arsitektur perangkat lunak beserta aspek yang
dinilainya, lalu menutup dengan bentuk pertama yang dibahas dalam modul ini,
yaitu client-server. Bentuk tersebut dipilih lebih dahulu karena hampir seluruh
bentuk lain pada bab-bab berikutnya tetap memakainya sebagai dasar. Pembahasan
mengikuti kerangka yang disusun oleh Richards dan
Ford~\cite{richards2020fundamentals}.

\section{Aspek yang Dinilai pada Sebuah Arsitektur}

Sebuah arsitektur tidak dinilai benar atau salah, melainkan sesuai atau tidak
sesuai dengan kebutuhan. Penilaiannya memakai sejumlah aspek, dan aspek-aspek
tersebut sering saling bertentangan. Tabel~\ref{tab:aspek} merangkum lima aspek
yang paling sering dipakai.

\begin{table}[htbp]
  \centering
  \caption{Lima aspek yang paling sering dipakai untuk menilai arsitektur}
  \label{tab:aspek}
  \begin{tabularx}{\textwidth}{|l|X|X|}
    \hline
    \textbf{Aspek} & \textbf{Pertanyaannya} & \textbf{Diukur dengan} \\
    \hline
    Kemudahan perawatan & Berapa banyak berkas berubah untuk satu permintaan perubahan? & Jumlah berkas dan baris \\
    \hline
    Kemampuan skala & Apa yang terjadi ketika jumlah pengguna naik sepuluh kali? & Latensi pada beban bertingkat \\
    \hline
    Kinerja & Berapa lama satu permintaan dilayani? & Median dan persentil ke-95 \\
    \hline
    Ketersediaan & Apa yang gagal ketika satu bagian mati? & Jumlah proses yang terhenti \\
    \hline
    Kemudahan pengujian & Berapa besar bagian yang dapat diuji tanpa infrastruktur? & Cakupan pengujian \\
    \hline
  \end{tabularx}
\end{table}

Aspek-aspek tersebut jarang dapat dipenuhi seluruhnya sekaligus. Menaikkan
kemampuan skala biasanya menurunkan kemudahan perawatan, sebab sistem harus
dipecah menjadi lebih banyak bagian. Menaikkan kinerja sering menurunkan
kemudahan pengujian, sebab jalan pintas mulai bermunculan. Pekerjaan seorang
arsitek karena itu bukan mencari susunan terbaik, melainkan memilih aspek mana
yang dikorbankan.

Setiap bab pada modul ini mengukur dua atau tiga aspek dari tabel tersebut, lalu
melaporkan angkanya. Cara itu dipakai agar perbandingan antar bentuk arsitektur
berdiri di atas angka, bukan di atas pendapat.

\section{Studi Kasus Singkat}

Tiga kejadian berikut menunjukkan bahwa keputusan arsitektural berakibat nyata,
bukan sekadar soal kerapian kode.

\textbf{Netflix.} Pada 2008 layanan Netflix berhenti cukup lama akibat kerusakan
pada basis data tunggalnya. Seluruh sistem saat itu berupa satu kesatuan, dan
kegagalan satu bagian menghentikan seluruhnya. Sesudah kejadian tersebut sistem
dipecah menjadi banyak layanan terpisah, sehingga kegagalan satu layanan tidak
lagi menghentikan yang lain.

\textbf{Healthcare.gov.} Pada peluncurannya tahun 2013 situs pendaftaran
asuransi kesehatan Amerika Serikat tersebut nyaris tidak dapat dipakai. Salah
satu sebabnya, sistem memanggil banyak sistem luar secara berurutan dan menunggu
seluruh jawabannya sebelum menampilkan halaman. Satu sistem luar yang lambat
membuat seluruh halaman ikut lambat.

\textbf{Sistem perbankan.} Sistem yang memproses transaksi secara berkelompok
pada akhir hari menunda pembaruan saldo sampai kelompok tersebut selesai
diproses. Perpindahan ke pemrosesan berbasis event membuat saldo diperbarui
segera setelah transaksi terjadi. Bab~9 membahas bentuk tersebut secara khusus.

Ketiganya memperlihatkan pola yang sama. Yang gagal bukan algoritmanya,
melainkan susunan bagian-bagiannya beserta cara mereka saling menunggu.

\section{Arsitektur Client-Server}

Client-server membagi sistem menjadi dua peran. Klien meminta, server melayani.
Gambar~\ref{fig:client-server} menunjukkan susunannya. Banyak klien memanggil
satu server, dan seluruh panah berawal dari klien.

\begin{figure}[htbp]
  \centering
  \begin{tikzpicture}[
    kotak/.style={draw, rounded corners=2pt, minimum width=1.8cm,
      minimum height=0.7cm, align=center, font=\scriptsize},
    server/.style={draw, rounded corners=3pt, minimum width=3.0cm,
      minimum height=1.4cm, align=center, font=\scriptsize, fill=praditagreen!15},
    panah/.style={-{Stealth[length=2mm]}, thick}
  ]
    \node[server] (srv) {\textbf{Server}\\
      {\tiny\mdseries menyimpan tabel kurs}\\
      {\tiny\mdseries menjalankan perhitungan}};

    \node[kotak, fill=blue!8, left=2.6cm of srv, yshift=1.15cm] (k1) {Klien 1};
    \node[kotak, fill=blue!8, left=2.6cm of srv] (k2) {Klien 2};
    \node[kotak, fill=blue!8, left=2.6cm of srv, yshift=-1.15cm] (k3) {Klien 3};

    \draw[panah] (k1.east) -- node[above, font=\tiny, pos=0.55] {permintaan} (srv.west);
    \draw[panah] (k2.east) -- (srv.west);
    \draw[panah] (k3.east) -- (srv.west);
  \end{tikzpicture}
  \caption{Banyak klien memanggil satu server. Seluruh permintaan berawal dari
    klien, dan server tidak pernah memanggil klien}
  \label{fig:client-server}
\end{figure}

Pembagian tersebut menentukan siapa menyimpan apa. Server menyimpan data
bersama dan menjalankan aturan bisnis. Klien menyimpan keadaan tampilan saja.
Karena tabel kurs hanya ada di server, memperbaruinya cukup dilakukan di satu
tempat, dan seluruh klien langsung memakai nilai baru tersebut.

Analogi yang mendekati adalah warung dengan satu kasir. Pembeli antre dan
memesan, kasir yang menghitung dan menyimpan uang. Pembeli tidak perlu tahu
harga pokok setiap barang, dan perubahan harga cukup diberitahukan kepada kasir.
Kelemahannya juga sama, yaitu ketika pembeli banyak, antrean memanjang walaupun
setiap pesanan sebenarnya cepat dilayani.

Listing~\ref{lst:penangan} memperlihatkan inti server pada contoh bab ini.
Server membaca parameter permintaan, menghitung, lalu membalas dalam bentuk
JavaScript Object Notation (JSON).

\begin{lstlisting}[style=PythonStyle, caption={Penanganan satu permintaan pada
  server, tersedia di \berkas{sources/chapter01/server.py}}, label=lst:penangan]
def do_GET(self):
  """Membaca parameter permintaan, menghitung, lalu membalas dalam JSON."""
  parameter = parse_qs(urlparse(self.path).query)
  asal = parameter.get("asal", ["USD"])[0]
  tujuan = parameter.get("tujuan", ["IDR"])[0]
  nominal = float(parameter.get("nominal", ["1"])[0])

  hasil = hitung_konversi(asal, tujuan, nominal)
  if hasil is None:
    self.kirim(404, {"galat": "Kurs tidak tersedia"})
  else:
    self.kirim(200, {"asal": asal, "tujuan": tujuan, "hasil": hasil})
\end{lstlisting}

Klien pada contoh tersebut sangat tipis. Klien hanya menyusun alamat, mengirim
permintaan, lalu menampilkan jawaban. Tidak ada tabel kurs dan tidak ada
perhitungan di sisi klien.

\section{Dua Biaya yang Melekat}

Pemisahan klien dan server menimbulkan dua biaya yang tidak dapat dihindari.
Keduanya menjadi bahan pengukuran pada latihan bab ini.

\textbf{Biaya perjalanan.} Setiap permintaan harus menempuh jaringan dua kali,
yaitu pergi dan pulang. Pada server di komputer yang sama biaya ini kecil, tetapi
tetap ada. Pada server yang berada di kota lain, biaya ini menjadi bagian
terbesar dari waktu tanggap, dan tidak dapat diperkecil dengan mempercepat
perhitungannya.

\textbf{Biaya antrean.} Satu server melayani banyak klien. Ketika permintaan
datang lebih cepat daripada kemampuan server melayaninya, sisanya menunggu.
Server pada contoh bab ini melayani satu permintaan pada satu waktu, sehingga
gejala tersebut muncul jelas begitu jumlah klien bertambah.

Tabel~\ref{tab:biaya-cs} merangkum hasil pengukuran nyata pada contoh bab ini.
Kolom terakhir memperlihatkan gejala antrean tersebut.

\begin{table}[htbp]
  \centering
  \caption{Latensi terhadap jumlah klien serentak pada server contoh}
  \label{tab:biaya-cs}
  \begin{tabularx}{\textwidth}{|X|l|l|}
    \hline
    \textbf{Jumlah klien serentak} & \textbf{Median (ms)} & \textbf{Terbesar (ms)} \\
    \hline
    1 & 0,240 & 8,027 \\
    \hline
    2 & 0,363 & 0,848 \\
    \hline
    4 & 0,777 & 1,748 \\
    \hline
    8 & 1,472 & 1011,138 \\
    \hline
    16 & 1,481 & 1019,625 \\
    \hline
  \end{tabularx}
\end{table}

Angka pada kolom terakhir melonjak tajam mulai delapan klien. Nilai median hanya
naik dari 0,24 menjadi 1,48 milidetik, sedangkan nilai terbesarnya mencapai
seribu milidetik. Sebagian permintaan karena itu menunggu sangat lama, sementara
sebagian besar lainnya tetap cepat. Gejala tersebut tidak terlihat bila yang
dilaporkan hanya nilai rata-rata, dan itulah alasan pengukuran pada modul ini
selalu menyertakan nilai terbesar atau persentil ke-95.

\section{Kelebihan dan Kekurangan}

Tabel~\ref{tab:untung-rugi-cs} merangkum untung rugi bentuk client-server.

\begin{table}[htbp]
  \centering
  \caption{Kelebihan dan kekurangan arsitektur client-server}
  \label{tab:untung-rugi-cs}
  \begin{tabularx}{\textwidth}{|X|X|}
    \hline
    \textbf{Kelebihan} & \textbf{Kekurangan} \\
    \hline
    Data bersama disimpan di satu tempat, sehingga pembaruan cukup sekali & Server menjadi satu-satunya titik yang bila mati menghentikan seluruh layanan \\
    \hline
    Aturan bisnis tidak dapat diubah dari sisi klien & Setiap permintaan menanggung biaya perjalanan jaringan \\
    \hline
    Klien menjadi ringan dan mudah diganti & Server menjadi antrean bersama ketika klien bertambah \\
    \hline
    Perangkat klien tidak perlu bertenaga besar & Menambah kemampuan server lebih mahal daripada menambah klien \\
    \hline
  \end{tabularx}
\end{table}

\section{Ringkasan}

Arsitektur adalah susunan bagian besar sistem beserta aturan hubungan di
antaranya. Yang membedakannya dari rancangan rinci adalah biaya perubahan.
Keputusan yang murah diubah bukan keputusan arsitektural, sedangkan keputusan
yang menyentuh hampir seluruh berkas termasuk di dalamnya.

Sebuah arsitektur dinilai memakai beberapa aspek yang saling bertentangan, yaitu
kemudahan perawatan, kemampuan skala, kinerja, ketersediaan, dan kemudahan
pengujian. Aspek-aspek tersebut jarang dapat dipenuhi seluruhnya sekaligus,
sehingga pekerjaan seorang arsitek adalah memilih aspek mana yang dikorbankan.

Ketiga studi kasus pada bab ini memperlihatkan pola yang sama. Yang gagal bukan
algoritmanya, melainkan susunan bagian-bagiannya beserta cara mereka saling
menunggu. Perbaikannya pun bukan mempercepat perhitungan, melainkan mengubah
susunannya.

Client-server membagi sistem menjadi peminta dan pelayan. Server menyimpan data
bersama dan menjalankan aturan bisnis, sedangkan klien menyimpan keadaan
tampilan saja. Pembagian tersebut membuat pembaruan data cukup dilakukan di satu
tempat, dan membuat aturan bisnis tidak dapat diubah dari sisi klien.

Pembagian itu menimbulkan dua biaya tetap, yaitu perjalanan jaringan dan antrean
di server. Pengukuran pada bab ini menunjukkan biaya antrean tidak terlihat pada
nilai median, tetapi tampak jelas pada nilai terbesar. Karena itu setiap
pengukuran pada modul ini melaporkan sebaran, bukan satu angka rata-rata.

\section{Latihan}

Ketiga latihan berikut memakai satu basis kode yang sama di
\berkas{sources/chapter01/}, tetapi menyasar tiga masalah yang berbeda.
Latihan~1 memeriksa pemisahan tanggung jawab. Latihan~2 mengukur biaya
perjalanan. Latihan~3 mengukur perilaku server ketika klien bertambah. Seluruh
pengukuran hanya dijalankan terhadap server milik sendiri, tidak terhadap
sistem milik pihak lain. Hasil setiap mahasiswa dapat berbeda, dan yang
dilaporkan adalah pola perbandingannya.

Server dijalankan lebih dahulu pada satu terminal seperti pada
Listing~\ref{lst:jalankan-server}, lalu skrip lainnya pada terminal kedua.

\begin{lstlisting}[language=bash, caption={Menjalankan server pada terminal pertama},
  label=lst:jalankan-server]
cd sources/chapter01
python server.py
\end{lstlisting}

\subsection*{Latihan 1: Memetakan Tanggung Jawab}

Latihan ini membuktikan bahwa klien tidak menyimpan maupun menghitung apa pun.
Langkah kerjanya sebagai berikut.

\begin{enumerate}
  \item Jalankan \texttt{python klien.py USD IDR 100} pada terminal kedua, lalu
    catat jawabannya.
  \item Buka \berkas{sources/chapter01/klien.py}, lalu cari tabel kurs di
    dalamnya.
  \item Ubah satu nilai pada tabel kurs di \berkas{sources/chapter01/server.py},
    jalankan ulang server, lalu jalankan klien tanpa mengubah apa pun.
  \item Matikan server, lalu jalankan klien sekali lagi dan catat yang terjadi.
\end{enumerate}

Tabel~\ref{tab:lat1cs-contoh} memperlihatkan bentuk pengisian yang diharapkan.

\begin{table}[htbp]
  \centering
  \caption{Contoh hasil Latihan 1 yang sudah terisi}
  \label{tab:lat1cs-contoh}
  \begin{tabularx}{\textwidth}{|X|l|X|}
    \hline
    \textbf{Pertanyaan} & \textbf{Jawaban} & \textbf{Bukti} \\
    \hline
    Klien menyimpan tabel kurs & tidak & Tidak ada kamus KURS di \texttt{klien.py} \\
    \hline
    Klien ikut berubah setelah kurs diubah & tidak & Jawaban berubah tanpa menyunting klien \\
    \hline
  \end{tabularx}
\end{table}

Isikan hasil pengamatan pada Tabel~\ref{tab:lat1cs-hasil}.

\begin{table}[htbp]
  \centering
  \caption{Lembar isian Latihan 1}
  \label{tab:lat1cs-hasil}
  \begin{tabularx}{\textwidth}{|X|l|X|}
    \hline
    \textbf{Pertanyaan} & \textbf{Jawaban} & \textbf{Bukti} \\
    \hline
    Klien menyimpan tabel kurs & \dots & \dots \\
    \hline
    Klien menghitung konversi & \dots & \dots \\
    \hline
    Klien ikut berubah setelah kurs diubah & \dots & \dots \\
    \hline
    Yang terjadi ketika server dimatikan & \multicolumn{2}{X|}{\dots} \\
    \hline
  \end{tabularx}
\end{table}

Jawab dua pertanyaan berikut. Pertama, jelaskan keuntungan menyimpan tabel kurs
hanya di server. Kedua, jelaskan kerugiannya berdasarkan yang terjadi ketika
server dimatikan.

\subsection*{Latihan 2: Mengukur Biaya Perjalanan}

Latihan ini mengukur waktu bolak-balik satu permintaan. Langkah kerjanya sebagai
berikut.

\begin{enumerate}
  \item Jalankan \texttt{python ukur\_latensi.py}, lalu catat kelima angka yang
    dilaporkan.
  \item Ulangi pengukuran sebanyak tiga kali.
  \item Bandingkan nilai terkecil dengan nilai terbesar pada setiap pengukuran.
  \item Tambahkan perhitungan berulang di dalam \texttt{hitung\_konversi} agar
    server bekerja lebih lama, lalu ukur kembali.
\end{enumerate}

Tabel~\ref{tab:lat2cs-contoh} memperlihatkan hasil satu kali pengukuran nyata.

\begin{table}[htbp]
  \centering
  \caption{Contoh hasil Latihan 2 yang sudah terisi}
  \label{tab:lat2cs-contoh}
  \begin{tabularx}{\textwidth}{|X|l|l|l|}
    \hline
    \textbf{Pengukuran} & \textbf{Median} & \textbf{Persentil ke-95} & \textbf{Terbesar} \\
    \hline
    Ke-1, 200 permintaan & 0,237 ms & 0,351 ms & 7,391 ms \\
    \hline
  \end{tabularx}
\end{table}

Isikan hasil pengukuran pada Tabel~\ref{tab:lat2cs-hasil}.

\begin{table}[htbp]
  \centering
  \caption{Lembar isian Latihan 2}
  \label{tab:lat2cs-hasil}
  \begin{tabularx}{\textwidth}{|X|l|l|l|}
    \hline
    \textbf{Pengukuran} & \textbf{Median} & \textbf{Persentil ke-95} & \textbf{Terbesar} \\
    \hline
    Ke-1 & \dots & \dots & \dots \\
    \hline
    Ke-2 & \dots & \dots & \dots \\
    \hline
    Ke-3 & \dots & \dots & \dots \\
    \hline
    Setelah server dibuat lebih lambat & \dots & \dots & \dots \\
    \hline
  \end{tabularx}
\end{table}

Jawab dua pertanyaan berikut dengan menunjuk angka pada tabel. Pertama, nilai
terbesar pada pengukuran pertama jauh melebihi mediannya, sehingga jelaskan
sebabnya. Kedua, jelaskan bagian mana dari waktu tersebut yang tetap ada
walaupun perhitungan di server dipercepat.

\subsection*{Latihan 3: Menambah Jumlah Klien}

Latihan ini membuktikan bahwa server adalah antrean bersama. Langkah kerjanya
sebagai berikut.

\begin{enumerate}
  \item Jalankan \texttt{python ukur\_beban.py}, lalu catat median dan nilai
    terbesar pada setiap jumlah klien.
  \item Tandai jumlah klien saat nilai terbesar mulai melonjak.
  \item Bandingkan pertumbuhan median terhadap pertumbuhan nilai terbesar.
  \item Jelaskan mengapa keduanya tumbuh dengan pola yang berbeda.
\end{enumerate}

Tabel~\ref{tab:lat3cs-contoh} memperlihatkan hasil satu kali pengukuran nyata.

\begin{table}[htbp]
  \centering
  \caption{Contoh hasil Latihan 3 yang sudah terisi}
  \label{tab:lat3cs-contoh}
  \begin{tabularx}{\textwidth}{|X|l|l|}
    \hline
    \textbf{Jumlah klien} & \textbf{Median (ms)} & \textbf{Terbesar (ms)} \\
    \hline
    4 & 0,777 & 1,748 \\
    \hline
    8 & 1,472 & 1011,138 \\
    \hline
  \end{tabularx}
\end{table}

Isikan hasil pengukuran pada Tabel~\ref{tab:lat3cs-hasil}.

\begin{table}[htbp]
  \centering
  \caption{Lembar isian Latihan 3}
  \label{tab:lat3cs-hasil}
  \begin{tabularx}{\textwidth}{|X|l|l|}
    \hline
    \textbf{Jumlah klien} & \textbf{Median (ms)} & \textbf{Terbesar (ms)} \\
    \hline
    1 & \dots & \dots \\
    \hline
    2 & \dots & \dots \\
    \hline
    4 & \dots & \dots \\
    \hline
    8 & \dots & \dots \\
    \hline
    16 & \dots & \dots \\
    \hline
    Jumlah klien saat lonjakan mulai & \multicolumn{2}{l|}{\dots} \\
    \hline
  \end{tabularx}
\end{table}

Jawab tiga pertanyaan berikut dengan menunjuk angka pada tabel. Pertama, median
tumbuh perlahan sedangkan nilai terbesar melonjak, sehingga jelaskan apa yang
dialami sebagian kecil permintaan tersebut. Kedua, jelaskan mengapa gejala itu
hilang bila hanya nilai rata-rata yang dilaporkan. Ketiga, dengan menggabungkan
hasil ketiga latihan, sebutkan aspek dari Tabel~\ref{tab:aspek} yang paling
tertekan pada arsitektur client-server beserta angka yang menjadi dasarnya.
